\section{Retrofitting a production codebase: Tutors}\label{sec:tutors}

The receivers of Section~\ref{sec:receivers} were built with the ladder in mind. This section asks
whether the ladder can be climbed by a codebase that was not: an open-source teaching platform in
production, with years of history and few tests, into which coding agents were already
contributing. It reports what was retrofitted, in what order, what it reached and what it could
not, and how its release gate was extracted into a separate, reusable repository. Every figure
below is read from the two public repositories, the monorepo at commit \texttt{865d02f}
\citep{tutorsmonorepo} and the release harness at version 1.15.0, commit \texttt{c324525}
\citep{tutorsreleaseharness}. Where a figure is our own count rather than one the repositories
record, we say so.

\subsection{Context}

Tutors is a course reader and learning platform. Lecturers write courses in Markdown; a
generator turns a course folder into one \repo{tutors.json}; four SvelteKit applications
(\repo{reader}, \repo{catalogue}, \repo{live} and \repo{time}) render it, track presence and
report learning time. The stack is TypeScript, Svelte~5 and SvelteKit, six Deno packages
published to JSR, pnpm workspaces and Supabase (\repo{README.md}, \repo{ARCHITECTURE.md}). The
monorepo holds the four applications and 24 packages. By our count, product code in
\repo{apps/} and \repo{packages/} is about 24,500 lines of TypeScript and Svelte, excluding
tests, fixtures, declaration files and a vendored PDF worker; test code is about 33,500 lines.

The monorepo was founded on 27 July 2026 by consolidating several repositories; the
original reader repository has history from 2022 (\path{guides/evolution/tutors-revolution.md}). On 1 June 2026 it had no tests, no CI workflow and no Dockerfile, and
the whole ecosystem had 27 tests (\path{guides/testing/journey.md}). Coding agents took part in
the retrofit throughout: by our count of git trailers, 329 of the monorepo's 483 non-merge
commits, and all 144 of the harness's, were authored or co-authored by an agent. The
retrofit is therefore a second known use of the paper's setting, not a laboratory one: a
domain protected while agents change it.

\subsection{The retrofit, rung by rung}

Table~\ref{tab:tutors-binding} maps each rung pattern to the artefact that binds it in Tutors.
The order in which they arrived says something about what each one costs. The first moves, in early August, were
to replace call-sequence mocks with real implementations and a data-driven stand-in carrying 27
validation tests of its own (commit \texttt{786f1fe}), and to clear 138 type errors
(\texttt{95064ea}). These are the cheapest rungs: they make existing tests honest before new
ones are added. Structure came next. On 16 September a ``testing runway'' added architecture
rules, suite-health lints and completeness checks as tiers A to O, and one aggregate CI job,
\repo{ci-success}, as the required status (\texttt{c57b9eb}, \texttt{3bf806b}). Requirements
came last. Feature files had existed since August, but on 17 September all 24 were found to be
documentation only; binding them to product code (\texttt{adffc55}) found bugs in search, sort
and logging (\path{guides/evolution/tutors-revolution.md}, ``Spec-first''). EARS \texttt{Rule:}
blocks with permanent ids and an audit followed on 21 September, and on 26 September coverage and
mutation floors.

\begin{table*}[t]
  \centering\small
  \caption{The rung patterns as bound in Tutors, with the reach of each. Paths are in the monorepo
  unless marked \emph{harness}.}
  \label{tab:tutors-binding}
  \begin{tabularx}{\linewidth}{>{\raggedright\arraybackslash}p{2.7cm}|X|>{\raggedright\arraybackslash}p{4.1cm}}
    \hline
    \textbf{Rung pattern} & \textbf{Tutors artefact} & \textbf{Reach} \\
    \hline\hline
    \pattern{Worked Slice} (L1) & \path{docs/COURSE-PAGE-WALKTHROUGH.md}: one URL followed through every file to the rendered cards & a slice in prose, not a small runnable exemplar \\
    \hline
    \pattern{Short Manifest} (L2) & \repo{CONTRIBUTING.md} (2,043 words), \repo{ARCHITECTURE.md}; the \repo{ears-gherkin-dev} agent skill (962 words) & no \repo{AGENTS.md} or \repo{CLAUDE.md}; the skill loads on demand, not every session \\
    \hline
    \pattern{Executable Rule} (L3) & \repo{.dependency-cruiser.cjs}: ten rules (six layer rules, no app-to-app, no relative cross-workspace import, primitives off data packages, no cross-package cycle); knip; manifest parity & new code at L3; three older violations held in \repo{known-violations.txt} \\
    \hline
    \pattern{Closure Check} (L3) & EARS audit, \path{scripts/checks/ears-audit.ts}: every Rule has a scenario, no scenario outside a Rule, every Rule bound, every UI scenario has a Playwright test and the reverse; every claim cites a Rule that exists & 95 Rules, 193 scenarios in 34 feature files (our count); audit required in \repo{ci-success} since 23 Sep \\
    \hline
    \pattern{Honest Double} (L3) & \repo{MockSupabaseClient} with its own tests; \emph{harness}: a second, in-memory backend behind each seam, driving the same collector & data-driven, not one contract run on both sides \\
    \hline
    \pattern{Teaching Failure} (L4) & a \repo{comment} on every dependency rule; ratchet messages (``Fix them rather than adding them to \ldots''); an audit code table with the fix per code & L4 wherever a check exists \\
    \hline
    \pattern{Unrepresentable Mistake} (L5) & strict TypeScript and blocking \repo{svelte-check}; Zod validation of the environment & weak: \repo{loTypes} is typed \repo{string[]}, so \repo{LoType} is \repo{string} \\
    \hline
    \pattern{Single Gate} & locally \repo{pnpm lint}, \repo{pnpm test}, \repo{pnpm check}; in CI one required job over eleven; at release, the harness Gate & three commands, then one status \\
    \hline
    \pattern{Human Boundary} & CODEOWNERS on \repo{release/claims.yaml} and \repo{release/SOP.md}; \emph{harness}: CODEOWNERS on masks, mutants, \repo{gate.ts} and the contract; a broad claim needs a named person & declared; review not yet practised \\
    \hline
  \end{tabularx}
\end{table*}

\paragraph{Structure} The architecture rules state Tutors' layering (foundation, core, feature,
three UI layers, apps) as one list in \repo{.dependency-cruiser.cjs}, from which one rule per
layer is generated. Each rule carries a \repo{comment} that says why and what to do instead, and
each has a negative fixture under \path{tests/architecture/fixtures/}, a small workspace that
breaks it on purpose. The test asserts that every fixture is caught, that no clean file is
flagged, and that the live repository adds no violation beyond a baseline
(\path{tests/architecture/dependency-rules.test.ts}). Tutors is layered, like the receivers, and
enforces the layering by the same move: an \pattern{Executable Rule} rather than ports and
adapters.

\paragraph{Behaviour and requirements} Behaviour and requirements are carried together. A
requirement is the title of a Gherkin \texttt{Rule:}, written in EARS \citep{mavin2009ears} with
exactly one ``shall'', a named system and a permanent id (\repo{@rule-0110}); its scenarios are the
proof, bound to product code by \repo{vitest-cucumber}. Listing~\ref{lst:tutors-rule} shows one.
The audit is a \pattern{Closure Check} in both directions, and names each failure with a code and
its fix (\path{.claude/skills/ears-gherkin-dev/references/audit-codes.md}). It reaches past the
repository: a release claim cites a Rule id, and \repo{pnpm check:release-claims} fails when no
feature defines it (\path{guides/testing/EARS-METHODOLOGY.md}). A requirement is thus traced from
its sentence to its scenario, its changelog line and the observable difference it licenses.
Unlike the receivers, Tutors files Rules by persona rather than by layer, so it does not use the
EARS pattern as a placement oracle (Table~\ref{tab:ears}).

\begin{lstlisting}[caption={An EARS Rule in Tutors, from the test-quality ratchets feature file. The audit checks its form and tag; its scenarios check its truth. Steps are wrapped for width.},label=lst:tutors-rule]
@rule-0111 @ears-unwanted
Rule: If total or per-package test coverage falls below
  its recorded floor, then tutors shall fail the coverage
  check and name the scope and metric.
  Scenario: A package below its floor fails the check
    Given recorded coverage floors where
      "packages/jsr/model/src/**" holds statements at 83 percent
    When the measured statement coverage of
      "packages/jsr/model/src" is 82.9 percent
    Then the coverage check shall report "below-floor:
      packages/jsr/model/src/** statements 82.9% < floor 83%"
\end{lstlisting}

\paragraph{Climbing in steps} Two rules show the ladder being climbed deliberately rather than
all at once. Type checks ran with \repo{continue-on-error} until the type errors were cleared, and
became blocking on 18 September (\texttt{1bea114}). The EARS audit was added as a non-blocking
job on 21 September (\texttt{6de1236}) and joined the required \repo{ci-success} job on
23 September (\texttt{4600835}). In both cases the rule was declared and measured before it was
allowed to stop the line.

\subsection{What a retrofit adds: ratchets and self-tests}

A greenfield receiver can make a rule true everywhere on the day it is written. A retrofit
cannot. Tutors' answer to that is the part of the retrofit we would most like other teams to
copy. Nine baseline files list what
is known to be wrong when a check is introduced. The check fails on any violation not in the
baseline and on any baseline entry that no longer occurs, so the list can only shrink
(\path{scripts/checks/lib/ratchet.ts}). The nine held 169 entries on 21 September and 167 on
28 September (\repo{journey.md}). Coverage and mutation floors ratchet the other way: each is the
measured value rounded down, and a floor two or more points below the measurement fails as stale
(\path{tests/suite-health/coverage-floors.json}, \path{tests/mutation/mutation-floors.json}).
A ratchet lets a rule sit at L3 for new code on the day it arrives while old code is paid down
from a ledger. The receivers never needed this, and a codebase with history cannot do without it.

The second addition answers the paper's own evaluation method. Section~\ref{sec:evaluation}
seeded mistakes from outside the receivers to show the checks can fail. Tutors keeps that
evidence inside the repository. ``No tier without a negative fixture'' is a stated rule
(\repo{journey.md}). Mutation testing covers twelve targeted modules, breaking below 90\% with
a measured score of 93.6\%, and a nightly run over every library module, whose first baseline
scored 58.5\% over 95 modules (\path{guides/testing/MUTATION-TESTING.md}). The harness applies
the same rule to itself: a new engine does not merge without an A/A test, a planted change it
catches and a change it must not flag (\emph{harness} \path{docs/user-guide/09-extending.md}).

\subsection{Abstraction: the release harness}

The last rung Tutors added sits outside the monorepo (Figure~\ref{fig:tutors-harness}). Its
precursor was an in-repository check that the generator's output matched production's, byte for
byte (\path{tests/release/RELEASE-TESTING.md}). The release harness generalises that check to
everything a user can observe. It runs a candidate's images beside production's, drives the same
scripted journeys through both, captures nineteen kinds of artefact (DOM, screenshots, network,
console, headers, accessibility, focus order, timing, metrics, logs, persisted writes, SBOM,
vulnerabilities, container posture and others), masks known noise, and fails unless every
difference is covered by a claim (\emph{harness} \repo{README.md}). A claim names an artefact, a
scope and the Rule or changelog entry that justifies it (\repo{release/claims.yaml}). This is a
\pattern{Closure Check} over the release: every observable difference has a claim, and a claim
that matches nothing is reported as stale.

\begin{figure*}[t]
  \centering
  \resizebox{0.92\linewidth}{!}{% Tutors: the rungs kept inside the monorepo, and the judge extracted into its own repository.
% Colours from assertionladder.sty. Included by sections/06a-tutors.tex.
\begin{tikzpicture}[
    font=\footnotesize,
    box/.style={draw=#1, line width=0.9pt, fill=#1!7, rounded corners=2pt, text width=3.3cm,
                align=left, inner sep=4pt},
    repo/.style={draw=#1, line width=1pt, rounded corners=4pt, inner sep=6pt},
    arr/.style={-{Stealth[length=2.3mm]}, line width=0.8pt, #1},
    lbl/.style={font=\scriptsize\bfseries, anchor=south west, inner sep=1pt},
  ]
  % Inside the monorepo
  \node[box=structure] (arch) {\textbf{Structure}\\dependency rules, knip,\\blocking type checks};
  \node[box=behaviour, below=0.25cm of arch] (bdd) {\textbf{Behaviour}\\Gherkin scenarios bound\\to product code};
  \node[box=requirement, below=0.25cm of bdd] (ears) {\textbf{Requirements}\\EARS \texttt{Rule:} with ids;\\the EARS audit};
  \node[box=muted, below=0.25cm of ears] (rat) {\textbf{Ratchets}\\shrink-only baselines;\\coverage, mutation floors};
  \node[repo=jotblue, fit=(arch)(bdd)(ears)(rat), label={[lbl, text=jotblue]north west:{\repo{tutors-mono-repo}: rungs checked on every PR}}] (mono) {};

  % The interface
  \node[box=requirement, right=1.25cm of bdd, text width=2.9cm, align=center] (claims)
       {\repo{release/claims.yaml}\\\repo{rules.json}\\signed images};

  % The harness
  \node[box=behaviour, right=1.25cm of claims, yshift=0.95cm] (gate) {\textbf{Gate}\\every observable\\difference claimed};
  \node[box=missed, below=0.25cm of gate] (mut) {\textbf{Self-test}\\ten planted mutants;\\a clean A/A first};
  \node[box=muted, below=0.25cm of mut] (kai) {\textbf{Kaizen}\\SOP, glance, 5 Whys,\\A3, scoreboard};
  \node[repo=missed, fit=(gate)(mut)(kai), label={[lbl, text=missed]north west:{\repo{tutors-release-harness}: the judge, kept apart}}] (har) {};

  \draw[arr=requirement] (ears.east) -- ++(0.35,0) |- (claims.west);
  \draw[arr=requirement] (claims.east) -- ++(0.35,0) |- (gate.west);
  \draw[arr=muted, dashed] (kai.south) |- node[pos=0.72, below, font=\scriptsize, align=center]{a 5 Whys ends in a countermeasure:\\a mutant, an EARS spec, an SOP change\ldots} (rat.east);
\end{tikzpicture}
}
  \caption{Tutors keeps structure, behaviour and requirement checks in the monorepo and ratchets
  their debt. The judge of a release is a separate repository that reads only signed images, the
  claims and the Rule ids, and whose 5 Whys feed countermeasures back into both.}
  \label{fig:tutors-harness}
\end{figure*}

Why a separate repository? The reason the harness gives is the one we gave in
Section~\ref{sec:ladder}: a check that an agent can edit is a check that an agent can weaken. The harness ``consumes the monorepo's
published images \ldots\ and knows nothing about their source, so it can compare any two tags
\ldots\ and cannot be quietly weakened by the PR it is judging'' (\emph{harness} \repo{README.md};
decision 1 in \path{docs/architecture/06-decisions.md}). The receivers bind the
\pattern{Human Boundary} at L2, as a line in the manifest. Extraction binds part of it
structurally: the gate's code, its masks and its mutants live where the change under review
cannot reach them, behind their own CODEOWNERS. The harness also holds itself to the discipline
it imposes. It may fail a release only after a clean production-against-production run less than
seven days old, and it must catch ten planted regressions, each attributed to the expected
artefact, weekly and on any change to an engine (\emph{harness} \repo{mutants/mutants.yaml},
\repo{src/gate.ts}).

What is generic is the contract, not the stack. The integration contract is versioned and
``is what \path{tutors-sdk/tutors-mono-repo} (or anything else) may build against''; a test fails
when document and code drift (\emph{harness} \repo{docs/contract.md}). An adopter supplies signed
images, a claims file and optionally a \repo{rules.json} of Rule ids, and triggers the harness by
dispatch (\emph{harness} \path{docs/monorepo/README.md}). The four application names, the
journeys and the fixture stubs are still Tutors'. We therefore read the harness as the paradigm
made reusable at the level of its contract, and not yet as a tool demonstrated on a second
product.

\paragraph{Lean, named in the code} The harness is explicit about its lean vocabulary, and each
term points to a mechanism rather than sitting in a glossary (\emph{harness} \repo{docs/lean.md};
\repo{release/SOP.md}). \emph{Jidoka}: the Gate stops the line, and an advisory confidence score
may never override it (``a FAIL is a FAIL at RCS 99''). \emph{Standard work}: a twelve-step
release SOP with an owner, an input and a done-when per step, in which steps 2 to 7 run as one
command that stops before step 8, the one step that must stay human. \emph{Gemba}: a reviewer's glance of
at most seven ranked places, each marked verified, disputed or escalated. \emph{Kaizen} and root
cause analysis: \repo{harness why} opens a 5 Whys whose first answer is the harness's own trace,
rejects ``human error'' as an answer, and requires one of seven countermeasure kinds, a new
mutant being preferred. \emph{Value stream}: \repo{harness a3} draws an A3 whose current condition
is a value stream map of a change from merge to production. Compression is written in as well:
``six journeys in three sets, not a hundred; ten mutants, not thirty. \ldots\ Add a journey only
when a real regression escaped that a journey would have caught'' (\emph{harness}
\repo{README.md}, ``Where to stop'').

\subsection{What the retrofit shows}

Three results are measured by the repositories themselves. First, the ladder can be climbed on a
production codebase in weeks: from 27 tests on 1 June to 2,140 declared tests, 184 bound scenarios
under 89 EARS Rules and 51 CI jobs on 28 September (\repo{journey.md}). Second, the release rung
finds what the lower rungs miss. The harness found four production bugs that no test had,
including an authentication page that answered 500 and a first visit that was never recorded
(\repo{journey.md}, ``Judge every release against production''). Third, it applies to a real
change the same root-cause discipline Section~\ref{sec:discussion} applies to the seeding misses:
three 5 Whys are open in the register, each ending in a countermeasure to the system rather than
to a person (\emph{harness} \repo{kaizen/README.md}).

Two results are demonstrated rather than measured. The three kinds of intent survive translation
to a codebase the ladder did not shape, and a requirement id can carry intent from a sentence
through a scenario to a release decision. And the retrofit adds two things the receivers lack
and a greenfield project would benefit from: ratchets for legacy debt, and self-tests that keep
the seeding evidence inside the repository.

We narrow the example to TypeScript because Tutors is written in it. The mechanisms have
counterparts in the other receivers' stacks: the dependency rules correspond to ArchUnit and to
the Python receiver's syntax-tree rules, and mutation tools exist for the JVM and for Python.
The harness judges images and claims, not source, so we expect a Java or Python service to adopt
it through the same contract; that is a hypothesis, not yet tried.

\subsection{Limits}

The retrofit does not reach every rung, and its own documents say so. There is no short
manifest loaded every session; the nearest are a 2,043-word contributing guide and one agent
skill loaded on demand. L5 is weak where the domain is richest: the learning-object type that
the whole model turns on resolves to \repo{string}
(\path{packages/jsr/model/src/types/type-utils.ts}). The global coverage floor was 58\% of
statements against a 90\% target on 28 September, and tier F, authorisation, is planned but not built
(\repo{journey.md}). The \pattern{Human Boundary} is declared but not yet practised: none of the
55 pull requests merged since release 16.2.2 had an approving review, and ``the CI gate, the
ratchets and the harness stand in for review; they do not replace it'' (\repo{journey.md}). The
release-confidence tooling is built, but its evidence gates need real releases: the scoreboard
is empty (\emph{harness} \path{scoreboard/releases.jsonl}), no 5 Whys is yet closed, and the daily
forecast of main against production was FAIL with 892 unclaimed differences on 28 September,
most of them an intended image change that nobody had claimed
(\path{guides/evolution/tutors-revolution.md}). Finally, Tutors is one codebase, the retrofit was
carried out largely by two of this paper's authors working with agents, and it was not designed as an
experiment. We report it as a known use of the pattern, not as a test of it.
